% source: RKF/theorum/42_cut_graded_lambda_jacobian_tower_theorem.md
\hypertarget{cut-graded-lambda-jacobian-tower-and-thermodynamic-recognition-theorem}{%
\chapter{Cut-Graded Lambda-Jacobian Tower and Thermodynamic Recognition Theorem}\label{cut-graded-lambda-jacobian-tower-and-thermodynamic-recognition-theorem}}

\hypertarget{source-problem-and-scope}{%
\section{1. Source problem and scope}\label{source-problem-and-scope}}

The thermodynamic paper \texttt{arXiv:2603.20773v2} introduces:

\begin{verbatim}
a signed T-V-S-P compass;
a four-component lambda response map;
a first Jacobian layer;
recursive higher derivative layers;
a jet-bundle analogy;
response one-forms and curvature-based cyclic residues.
\end{verbatim}

The source manuscript explicitly describes the jet-bundle relation as a
geometric analogy rather than a rigorous construction. It also treats the
nonequilibrium curvature and information-sector extensions as constitutive,
calibrated hypotheses requiring further validation.

The present theorem supplies the missing abstract object:

\begin{verbatim}
smooth state manifold and response bundle
-> declared cut on state and response fibres
-> covariant derivative raising operator
-> typed recursive response tower
-> induced cut grading at every derivative order
-> cut/transpose bi-grading of the Jacobian
-> cut-loop mixed curvature
-> target-relative finite tower observer.
\end{verbatim}

This theorem does \textbf{not} prove that the particular physical lambda map in the
source paper is globally antisymmetric, that its proposed curvature is already
entropy production, or that every phase transition is a rank-change cut.
Those statements require separate constitutive and experimental pins.

\begin{center}\rule{0.5\linewidth}{0.5pt}\end{center}

\hypertarget{native-thermodynamic-data}{%
\section{2. Native thermodynamic data}\label{native-thermodynamic-data}}

Let \texttt{M} be a smooth four-dimensional state manifold. On a declared chart write

\[
x=(T,V,S,P).
\]

The coordinates must be nondimensionalized or treated as sections of their
proper dimensional bundles before a linear involution is imposed.

Let \texttt{E\ -\textgreater{}\ M} be a rank-four response bundle and let

\[
\Lambda\in\Gamma(E)
\]

be a smooth response section. In the source application, the local components
are the four constraint-direction responses

\[
\Lambda=(\lambda_p,\lambda_v,\lambda_s,\lambda_t),
\]

up to the declared ordering convention.

Let

\[
j:M\longrightarrow M,
\qquad j^2=\operatorname{id}_M,
\]

be a state-space cut and let

\[
J_E:E\longrightarrow E,
\qquad J_E^2=I,
\]

cover \texttt{j}. Let \texttt{nabla} be a connection on \texttt{E} satisfying cut covariance:

\[
\boxed{
\nabla(j^*s)=j^*(\nabla s)
}
\tag{2.1}
\]

in the induced tensor representation.

A compass candidate is the local state involution

\[
J_x=\operatorname{diag}(1,1,-1,-1),
\tag{2.2}
\]

but the physical state and response involutions must be derived from the
signed thermodynamic axes and checked against the lambda definitions. They are
not licensed merely by visual symmetry.

\begin{center}\rule{0.5\linewidth}{0.5pt}\end{center}

\hypertarget{the-rigorous-recursive-tower}{%
\section{3. The rigorous recursive tower}\label{the-rigorous-recursive-tower}}

Define the response layers by

\[
L_1:=\Lambda,
\qquad
L_{n+1}:=\nabla L_n
\quad(n\ge1).
\tag{3.1}
\]

Then

\[
L_n\in
\Gamma\!\left(
E\otimes(T^*M)^{\otimes(n-1)}
\right).
\tag{3.2}
\]

Thus:

\begin{verbatim}
L1  response section;
L2  covariant Jacobian;
L3  covariant Hessian;
L4  third response derivative;
...
\end{verbatim}

\hypertarget{theorem-3.1-typed-lambda-jacobian-tower}{%
\subsection{Theorem 3.1 (Typed lambda-Jacobian tower)}\label{theorem-3.1-typed-lambda-jacobian-tower}}

For every smooth response section \texttt{Lambda}, equation (3.1) defines a unique
covariant derivative tower. In a flat chart,

\[
L_n=D^{n-1}\Lambda.
\tag{3.3}
\]

If \texttt{Lambda} is polynomial of degree \texttt{d} in that chart, then

\[
L_n=0
\qquad(n>d+1).
\tag{3.4}
\]

The formal exponential generating object is

\[
\boxed{
\mathcal T_t\Lambda
:=
\sum_{k=0}^{\infty}
\frac{t^k}{k!}\nabla^k\Lambda,
}
\tag{3.5}
\]

whenever the series converges in the chosen tensor topology. The recursive
layers are its Taylor coefficients. The tower is therefore generated by one
raising operator rather than by repeatedly reusing a diagrammatic rule without
a type map.

\hypertarget{proof}{%
\subsubsection{Proof}\label{proof}}

A connection extends uniquely to tensor products by the Leibniz rule. Repeated
application gives (3.2). In a flat chart the connection coefficients vanish,
so the layers are ordinary higher derivatives. Polynomial termination follows
from elementary differentiation. Equation (3.5) is the exponential series of
the raising operator on the completed direct sum of tensor levels. \texttt{square}

\begin{center}\rule{0.5\linewidth}{0.5pt}\end{center}

\hypertarget{induced-cut-on-every-layer}{%
\section{4. Induced cut on every layer}\label{induced-cut-on-every-layer}}

Let \texttt{J\_x\^{}*} denote the induced involution on covectors. On layer \texttt{n}, define

\[
\boxed{
\mathbb J_n
=J_E\otimes(J_x^*)^{\otimes(n-1)}.
}
\tag{4.1}
\]

Then

\[
\mathbb J_n^2=I.
\]

Define

\[
\boxed{
L_n^{\pm}
=
\frac12\left(L_n\pm\mathbb J_nL_n\right).
}
\tag{4.2}
\]

\hypertarget{theorem-4.1-cut-graded-response-tower}{%
\subsection{Theorem 4.1 (Cut-graded response tower)}\label{theorem-4.1-cut-graded-response-tower}}

Every response layer has the unique decomposition

\[
\boxed{L_n=L_n^++L_n^-.}
\tag{4.3}
\]

Moreover,

\[
\mathbb J_nL_n^+=L_n^+,
\qquad
\mathbb J_nL_n^-=-L_n^-.
\tag{4.4}
\]

If the response section is cut equivariant,

\[
\Lambda(jx)=J_E\Lambda(x),
\tag{4.5}
\]

then for every \texttt{k\ \textgreater{}=\ 1},

\[
D^k\Lambda(jx)
[J_xh_1,\ldots,J_xh_k]
=
J_E D^k\Lambda(x)[h_1,\ldots,h_k].
\tag{4.6}
\]

At a fixed seam \texttt{jx=x}, every derivative layer obeys the induced parity
selection rule.

\hypertarget{proof-1}{%
\subsubsection{Proof}\label{proof-1}}

Equations (4.2)--(4.4) are the spectral projections of an involution.
Differentiate (4.5) repeatedly and apply the chain rule to obtain (4.6).
\texttt{square}

Interpretation:

\begin{verbatim}
L_n^+   response derivatives preserving the declared cut class;
L_n^-   response derivatives transporting across the declared cut class.
\end{verbatim}

This grading is defined at every order, not only at the first Jacobian.

\begin{center}\rule{0.5\linewidth}{0.5pt}\end{center}

\hypertarget{the-jacobian-is-bi-graded-not-automatically-antisymmetric}{%
\section{5. The Jacobian is bi-graded, not automatically antisymmetric}\label{the-jacobian-is-bi-graded-not-automatically-antisymmetric}}

Assume locally that the state and response fibres have the same dimension and
use one declared involution \texttt{J}. Let

\[
A=D\Lambda(x)\in M_4(\mathbb R).
\]

Define two commuting involutions on matrices:

\[
\mathcal C(A)=JAJ,
\qquad
\mathcal T(A)=A^{\mathsf T}.
\tag{5.1}
\]

For \texttt{sigma,tau\ in\ \{+1,-1\}}, define

\[
\boxed{
A^{\sigma,\tau}
=
\frac14
\left(
A+\sigma JAJ
+\tau A^{\mathsf T}
+\sigma\tau JA^{\mathsf T}J
\right).
}
\tag{5.2}
\]

\hypertarget{theorem-5.1-cuttranspose-jacobian-decomposition}{%
\subsection{Theorem 5.1 (Cut/transpose Jacobian decomposition)}\label{theorem-5.1-cuttranspose-jacobian-decomposition}}

The four sectors satisfy

\[
\boxed{
A=\sum_{\sigma,\tau\in\{\pm1\}}
A^{\sigma,\tau},
}
\tag{5.3}
\]

\[
JA^{\sigma,\tau}J
=\sigma A^{\sigma,\tau},
\tag{5.4}
\]

and

\[
(A^{\sigma,\tau})^{\mathsf T}
=\tau A^{\sigma,\tau}.
\tag{5.5}
\]

Therefore the Jacobian contains four logically distinct pieces:

\begin{verbatim}
cut-even / symmetric;
cut-even / antisymmetric;
cut-odd  / symmetric;
cut-odd  / antisymmetric.
\end{verbatim}

Antisymmetry

\[
A^{\mathsf T}=-A
\tag{5.6}
\]

is equivalent to the vanishing of both symmetric sectors. It is a theorem to
be proved from the constitutive lambda map and coordinate convention, not an
automatic consequence of having a T-V-S-P compass.

\hypertarget{proof-2}{%
\subsubsection{Proof}\label{proof-2}}

The maps \texttt{C} and \texttt{T} commute and square to the identity. Equation (5.2) is the
joint spectral projection formula for two commuting involutions. Equations
(5.3)--(5.5) follow immediately. \texttt{square}

This theorem separates Maxwell-type antisymmetric response from symmetric
susceptibility or dissipative response and separately identifies whether each
piece is within-sheet or trans-cut.

\begin{center}\rule{0.5\linewidth}{0.5pt}\end{center}

\hypertarget{cut-loop-curvature-and-a-lawful-thermodynamic-adapter}{%
\section{6. Cut-loop curvature and a lawful thermodynamic adapter}\label{cut-loop-curvature-and-a-lawful-thermodynamic-adapter}}

Let \texttt{G} be a response-space protocol generator and decompose it using the
certified central theorem:

\[
G=G_{\mathrm e}+G_{\mathrm o},
\qquad
JG_{\mathrm e}J=G_{\mathrm e},
\qquad
JG_{\mathrm o}J=-G_{\mathrm o}.
\tag{6.1}
\]

Let

\[
U_t=e^{tG},
\qquad
\mathscr H_J(t)=JU_tJU_t.
\tag{6.2}
\]

Then

\[
\log\mathscr H_J(t)
=
2tG_{\mathrm e}
+t^2[G_{\mathrm e},G_{\mathrm o}]
+O(t^3).
\tag{6.3}
\]

Now consider the constant connection one-form on a two-parameter protocol
chart \texttt{(s,t)}:

\[
\mathcal A
=G_{\mathrm e}\,ds+G_{\mathrm o}\,dt.
\tag{6.4}
\]

Its curvature is

\[
\mathcal F
=d\mathcal A+\mathcal A\wedge\mathcal A,
\]

so

\[
\boxed{
\mathcal F_{st}
=[G_{\mathrm e},G_{\mathrm o}].
}
\tag{6.5}
\]

\hypertarget{theorem-6.1-cut-loopconnection-curvature-identity}{%
\subsection{Theorem 6.1 (Cut-loop/connection curvature identity)}\label{theorem-6.1-cut-loopconnection-curvature-identity}}

For the constant cut-graded connection (6.4), the quadratic coefficient in the
logarithmic cut loop is exactly the mixed curvature component:

\[
\boxed{
[t^2]\log\mathscr H_J(t)
=
\mathcal F_{st}.
}
\tag{6.6}
\]

\hypertarget{proof-3}{%
\subsubsection{Proof}\label{proof-3}}

Equation (6.3) is the cut-loop extraction theorem. For constant connection
coefficients, derivative terms vanish and the mixed curvature is the
commutator. \texttt{square}

This is an exact finite-dimensional curvature model. To identify it with a
thermodynamic response two-form such as

\[
\Omega_H=d\omega_H,
\]

one must construct and prove an adapter

\[
\boxed{
\rho\bigl([G_{\mathrm e},G_{\mathrm o}]\bigr)
=\Omega_H.
}
\tag{6.7}
\]

Equation (6.7) is the next physical theorem. It is not supplied merely by the
abstract commutator identity.

\begin{center}\rule{0.5\linewidth}{0.5pt}\end{center}

\hypertarget{finite-tower-observers-and-recognition-kernels}{%
\section{7. Finite tower observers and recognition kernels}\label{finite-tower-observers-and-recognition-kernels}}

Let \texttt{V} be a finite-dimensional model or perturbation space. For each tower
level let

\[
Q_nL_n:V\longrightarrow Y_n
\]

be the declared measurable response packet at order \texttt{n}. Define the order-\texttt{N}
observer

\[
\boxed{
\mathcal A_Nv
=
(Q_1L_1v,\ldots,Q_NL_Nv).
}
\tag{7.1}
\]

For a target map \texttt{Pi:V\ -\textgreater{}\ Z}, define the target-relative blind quotient

\[
\boxed{
\mathfrak B_N
=
\frac{\ker\mathcal A_N}
{\ker\mathcal A_N\cap\ker\Pi}.
}
\tag{7.2}
\]

\hypertarget{theorem-7.1-tower-monotonicity-and-target-repair}{%
\subsection{Theorem 7.1 (Tower monotonicity and target repair)}\label{theorem-7.1-tower-monotonicity-and-target-repair}}

The recognition kernels are nested:

\[
\boxed{
\ker\mathcal A_{N+1}
\subseteq
\ker\mathcal A_N.
}
\tag{7.3}
\]

A higher response layer repairs a first-order blind target precisely when

\[
\ker\mathcal A_{N+1}\cap\operatorname{supp}\Pi
\subsetneq
\ker\mathcal A_N\cap\operatorname{supp}\Pi.
\tag{7.4}
\]

For a scalar target \texttt{L}, the Recognition-Kernel burden is

\[
\beta_N(L)
=
\sup_{\mathcal A_Nv\ne0}
\frac{|L(v)|^2}{\|\mathcal A_Nv\|^2}.
\tag{7.5}
\]

Whenever it is finite, the minimum decoder exists and

\[
\mathcal A_N^*\mathcal A_N-L^*L\ge0
\iff
\beta_N(L)\le1.
\tag{7.6}
\]

\hypertarget{proof-4}{%
\subsubsection{Proof}\label{proof-4}}

Adding an observer component can only shrink the common kernel, proving (7.3).
Equation (7.4) is the definition of strict target-relevant repair. Equation
(7.6) is the Recognition-Kernel Theorem applied to the direct-sum tower
observer. \texttt{square}

This theorem asks experimentally meaningful questions:

\begin{verbatim}
Which derivative order first recognizes a target?
Which response components form the minimum faithful packet?
Does a higher-order susceptibility repair a first-order blind direction?
\end{verbatim}

\begin{center}\rule{0.5\linewidth}{0.5pt}\end{center}

\hypertarget{rank-change-cuts}{%
\section{8. Rank-change cuts}\label{rank-change-cuts}}

Let \texttt{A\_1(r)} be a first-layer observer depending continuously on a control
parameter \texttt{r}. A state \texttt{r\_c} is a first-layer rank-change cut when

\[
\operatorname{rank}A_1(r_c)
<
\limsup_{r\to r_c}\operatorname{rank}A_1(r).
\tag{8.1}
\]

Let \texttt{A\_\{\textbackslash{}le\ N\}(r)} be the observer obtained by stacking response layers through
order \texttt{N}.

\hypertarget{proposition-8.1-higher-layer-recovery-at-a-rank-change-cut}{%
\subsection{Proposition 8.1 (Higher-layer recovery at a rank-change cut)}\label{proposition-8.1-higher-layer-recovery-at-a-rank-change-cut}}

If

\[
\ker A_1(r_c)\cap\ker A_{\le N}(r_c)=\{0\},
\tag{8.2}
\]

then the higher-order tower resolves every first-order blind direction at the
cut.

\hypertarget{proof-5}{%
\subsubsection{Proof}\label{proof-5}}

A vector blind to the first layer but also blind to the full tower lies in the
intersection in (8.2), hence must vanish. \texttt{square}

A thermodynamic phase transition may motivate such a test because the source
paper restricts its smooth lambda map away from response singularities.
However, calling every phase transition a rank-change cut requires data and a
specific constitutive observer; it is not an abstract consequence of (8.1).

\begin{center}\rule{0.5\linewidth}{0.5pt}\end{center}

\hypertarget{exact-rational-certificate}{%
\section{9. Exact rational certificate}\label{exact-rational-certificate}}

The proof-lab packet uses a cut-equivariant quadratic response map on four
coordinates and verifies exactly:

\begin{verbatim}
value-level cut equivariance;
Jacobian covariance;
induced Hessian covariance;
typed recursive layers through Hessian order;
termination of the quadratic tower;
four-sector cut/transpose Jacobian reconstruction;
nonzero symmetric and antisymmetric sectors;
negative control rejecting automatic Jacobian antisymmetry;
constant-connection curvature equals cut-loop seam curvature;
first-order target blindness;
second-layer target repair;
minimum decoder burden 1/4 and reserve 3/4;
first-layer rank loss at a cut;
higher-layer recovery of the cut kernel.
\end{verbatim}

No NumPy, floating point, fitted coordinate transform, or post-hoc Gram factor
is used.

Expected status:

\begin{verbatim}
PASS_CUT_GRADED_LAMBDA_JACOBIAN_TOWER_CANDIDATE
\end{verbatim}

Expected SHA-256:

\begin{verbatim}
3d8db4f085a3b624192c45b9f6e0356b600834193967bc89478381b53ae44a3a
\end{verbatim}

\begin{center}\rule{0.5\linewidth}{0.5pt}\end{center}

\hypertarget{claim-boundary}{%
\section{10. Claim boundary}\label{claim-boundary}}

\begin{verbatim}
TYPED COVARIANT RESPONSE TOWER                         PROVED
INDUCED CUT GRADING AT EVERY ORDER                     PROVED
CUT/TRANSPOSE JACOBIAN BI-GRADING                      PROVED
ANTISYMMETRY AS A SECTOR CONDITION                     PROVED
CONSTANT-CONNECTION CUT-LOOP CURVATURE IDENTITY        PROVED
FINITE TOWER RECOGNITION-KERNEL THEOREM                PROVED
RANK-CHANGE / HIGHER-LAYER RECOVERY PROPOSITION        PROVED
EXACT RATIONAL CERTIFICATE                             IMPLEMENTED / USER RUN REQUIRED

PHYSICAL T-V-S-P CUT AND RESPONSE-FIBRE INVOLUTION     OPEN / NEXT
CUT-COMPATIBLE THERMODYNAMIC CONNECTION                OPEN / NEXT
OPERATOR CURVATURE -> RESPONSE TWO-FORM ADAPTER        OPEN / NEXT
GLOBAL ANTISYMMETRY OF THE PAPER'S PHYSICAL JACOBIAN   NOT PROMOTED
CURVATURE = ENTROPY PRODUCTION                         REQUIRES CALIBRATION
PHASE TRANSITION = RANK-CHANGE CUT                     TESTABLE, NOT CLAIMED
SOURCE PAPER UPDATE                                    HELD UNTIL CERTIFICATE PASS
\end{verbatim}

\begin{flushright}\small\texttt{RKF/theorum/42\_cut\_graded\_lambda\_jacobian\_tower\_theorem.md}\end{flushright}
